\documentclass[
  doc,
  floatsintext,
  longtable,
  a4paper,
  nolmodern,
  notxfonts,
  notimes,
  12pt,
  colorlinks=true,linkcolor=blue,citecolor=blue,urlcolor=blue]{apa7}

\usepackage{amsmath}
\usepackage{amssymb}

\geometry{inner=1in, outer=1in}
\fancyhfoffset[LE,RO]{0cm}


\usepackage[bidi=default]{babel}
\babelprovide[main,import]{spanish}


% get rid of language-specific shorthands (see #6817):
\let\LanguageShortHands\languageshorthands
\def\languageshorthands#1{}

\RequirePackage{longtable}
\RequirePackage{threeparttablex}

\makeatletter
\renewcommand{\paragraph}{\@startsection{paragraph}{4}{\parindent}%
	{0\baselineskip \@plus 0.2ex \@minus 0.2ex}%
	{-.5em}%
	{\normalfont\normalsize\bfseries\typesectitle}}

\renewcommand{\subparagraph}[1]{\@startsection{subparagraph}{5}{0.5em}%
	{0\baselineskip \@plus 0.2ex \@minus 0.2ex}%
	{-\z@\relax}%
	{\normalfont\normalsize\bfseries\itshape\hspace{\parindent}{#1}\textit{\addperi}}{\relax}}
\makeatother




\usepackage{longtable, booktabs, multirow, multicol, colortbl, hhline, caption, array, float, xpatch}
\usepackage{subcaption}


\renewcommand\thesubfigure{\Alph{subfigure}}
\setcounter{topnumber}{2}
\setcounter{bottomnumber}{2}
\setcounter{totalnumber}{4}
\renewcommand{\topfraction}{0.85}
\renewcommand{\bottomfraction}{0.85}
\renewcommand{\textfraction}{0.15}
\renewcommand{\floatpagefraction}{0.7}

\usepackage{tcolorbox}
\tcbuselibrary{listings,theorems, breakable, skins}
\usepackage{fontawesome5}

\definecolor{quarto-callout-color}{HTML}{909090}
\definecolor{quarto-callout-note-color}{HTML}{0758E5}
\definecolor{quarto-callout-important-color}{HTML}{CC1914}
\definecolor{quarto-callout-warning-color}{HTML}{EB9113}
\definecolor{quarto-callout-tip-color}{HTML}{00A047}
\definecolor{quarto-callout-caution-color}{HTML}{FC5300}
\definecolor{quarto-callout-color-frame}{HTML}{ACACAC}
\definecolor{quarto-callout-note-color-frame}{HTML}{4582EC}
\definecolor{quarto-callout-important-color-frame}{HTML}{D9534F}
\definecolor{quarto-callout-warning-color-frame}{HTML}{F0AD4E}
\definecolor{quarto-callout-tip-color-frame}{HTML}{02B875}
\definecolor{quarto-callout-caution-color-frame}{HTML}{FD7E14}

%\newlength\Oldarrayrulewidth
%\newlength\Oldtabcolsep


\usepackage{hyperref}



\usepackage{color}
\usepackage{fancyvrb}
\newcommand{\VerbBar}{|}
\newcommand{\VERB}{\Verb[commandchars=\\\{\}]}
\DefineVerbatimEnvironment{Highlighting}{Verbatim}{commandchars=\\\{\}}
% Add ',fontsize=\small' for more characters per line
\usepackage{framed}
\definecolor{shadecolor}{RGB}{241,243,245}
\newenvironment{Shaded}{\begin{snugshade}}{\end{snugshade}}
\newcommand{\AlertTok}[1]{\textcolor[rgb]{0.68,0.00,0.00}{#1}}
\newcommand{\AnnotationTok}[1]{\textcolor[rgb]{0.37,0.37,0.37}{#1}}
\newcommand{\AttributeTok}[1]{\textcolor[rgb]{0.40,0.45,0.13}{#1}}
\newcommand{\BaseNTok}[1]{\textcolor[rgb]{0.68,0.00,0.00}{#1}}
\newcommand{\BuiltInTok}[1]{\textcolor[rgb]{0.00,0.23,0.31}{#1}}
\newcommand{\CharTok}[1]{\textcolor[rgb]{0.13,0.47,0.30}{#1}}
\newcommand{\CommentTok}[1]{\textcolor[rgb]{0.37,0.37,0.37}{#1}}
\newcommand{\CommentVarTok}[1]{\textcolor[rgb]{0.37,0.37,0.37}{\textit{#1}}}
\newcommand{\ConstantTok}[1]{\textcolor[rgb]{0.56,0.35,0.01}{#1}}
\newcommand{\ControlFlowTok}[1]{\textcolor[rgb]{0.00,0.23,0.31}{\textbf{#1}}}
\newcommand{\DataTypeTok}[1]{\textcolor[rgb]{0.68,0.00,0.00}{#1}}
\newcommand{\DecValTok}[1]{\textcolor[rgb]{0.68,0.00,0.00}{#1}}
\newcommand{\DocumentationTok}[1]{\textcolor[rgb]{0.37,0.37,0.37}{\textit{#1}}}
\newcommand{\ErrorTok}[1]{\textcolor[rgb]{0.68,0.00,0.00}{#1}}
\newcommand{\ExtensionTok}[1]{\textcolor[rgb]{0.00,0.23,0.31}{#1}}
\newcommand{\FloatTok}[1]{\textcolor[rgb]{0.68,0.00,0.00}{#1}}
\newcommand{\FunctionTok}[1]{\textcolor[rgb]{0.28,0.35,0.67}{#1}}
\newcommand{\ImportTok}[1]{\textcolor[rgb]{0.00,0.46,0.62}{#1}}
\newcommand{\InformationTok}[1]{\textcolor[rgb]{0.37,0.37,0.37}{#1}}
\newcommand{\KeywordTok}[1]{\textcolor[rgb]{0.00,0.23,0.31}{\textbf{#1}}}
\newcommand{\NormalTok}[1]{\textcolor[rgb]{0.00,0.23,0.31}{#1}}
\newcommand{\OperatorTok}[1]{\textcolor[rgb]{0.37,0.37,0.37}{#1}}
\newcommand{\OtherTok}[1]{\textcolor[rgb]{0.00,0.23,0.31}{#1}}
\newcommand{\PreprocessorTok}[1]{\textcolor[rgb]{0.68,0.00,0.00}{#1}}
\newcommand{\RegionMarkerTok}[1]{\textcolor[rgb]{0.00,0.23,0.31}{#1}}
\newcommand{\SpecialCharTok}[1]{\textcolor[rgb]{0.37,0.37,0.37}{#1}}
\newcommand{\SpecialStringTok}[1]{\textcolor[rgb]{0.13,0.47,0.30}{#1}}
\newcommand{\StringTok}[1]{\textcolor[rgb]{0.13,0.47,0.30}{#1}}
\newcommand{\VariableTok}[1]{\textcolor[rgb]{0.07,0.07,0.07}{#1}}
\newcommand{\VerbatimStringTok}[1]{\textcolor[rgb]{0.13,0.47,0.30}{#1}}
\newcommand{\WarningTok}[1]{\textcolor[rgb]{0.37,0.37,0.37}{\textit{#1}}}

\providecommand{\tightlist}{%
  \setlength{\itemsep}{0pt}\setlength{\parskip}{0pt}}
\usepackage{longtable,booktabs,array}
\newcounter{none} % for unnumbered tables
\usepackage{calc} % for calculating minipage widths
% Correct order of tables after \paragraph or \subparagraph
\usepackage{etoolbox}
\makeatletter
\patchcmd\longtable{\par}{\if@noskipsec\mbox{}\fi\par}{}{}
\makeatother
% Allow footnotes in longtable head/foot
\IfFileExists{footnotehyper.sty}{\usepackage{footnotehyper}}{\usepackage{footnote}}
\makesavenoteenv{longtable}

\usepackage{graphicx}
\makeatletter
\newsavebox\pandoc@box
\newcommand*\pandocbounded[1]{% scales image to fit in text height/width
  \sbox\pandoc@box{#1}%
  \Gscale@div\@tempa{\textheight}{\dimexpr\ht\pandoc@box+\dp\pandoc@box\relax}%
  \Gscale@div\@tempb{\linewidth}{\wd\pandoc@box}%
  \ifdim\@tempb\p@<\@tempa\p@\let\@tempa\@tempb\fi% select the smaller of both
  \ifdim\@tempa\p@<\p@\scalebox{\@tempa}{\usebox\pandoc@box}%
  \else\usebox{\pandoc@box}%
  \fi%
}
% Set default figure placement to htbp
\def\fps@figure{htbp}
\makeatother







\usepackage{newtx}

\defaultfontfeatures{Scale=MatchLowercase}
\defaultfontfeatures[\rmfamily]{Ligatures=TeX,Scale=1}





\title{Modificación de cadenas}


\shorttitle{MODIFICACIÓN DE CADENAS}


\usepackage{etoolbox}



\ccoppy{\textcopyright~2026}



\author{Edison Achalma}



\affiliation{
{Escuela Profesional de Economía, Universidad Nacional de San Cristóbal
de Huamanga}}




\leftheader{Achalma}

\date{2026-10-20}


\abstract{Este abstract será actualizado una vez que se complete el
contenido final del artículo. }



\authornote{\par{\addORCIDlink{Edison Achalma}{0000-0001-6996-3364}} 
\par{ }
\par{   El autor no tiene conflictos de interés que revelar.    Los
roles de autor se clasificaron utilizando la taxonomía de roles de
colaborador (CRediT; https://credit.niso.org/) de la siguiente
manera:  Edison Achalma:   conceptualización, metodología, análisis
formal, investigación, recursos, curación de
datos, redacción, visualización, supervisión, administración del
proyecto}
\par{La correspondencia relativa a este artículo debe dirigirse a Edison
Achalma, Escuela Profesional de Economía, Universidad Nacional de San
Cristóbal de Huamanga, Portal Independencia N°
57, Ayacucho, AYA 5001, Perú, Email: \href{mailto:elmer.achalma.09@unsch.edu.pe}{elmer.achalma.09@unsch.edu.pe}}
}

\makeatletter
\let\endoldlt\endlongtable
\def\endlongtable{
\hline
\endoldlt
}
\makeatother

\urlstyle{same}



\hypersetup{
  pdftitle={Modificaci\'{o}n de cadenas},
  pdfauthor={Autor}
}
\makeatletter
\@ifpackageloaded{caption}{}{\usepackage{caption}}
\AtBeginDocument{%
\ifdefined\contentsname
  \renewcommand*\contentsname{Tabla de contenidos}
\else
  \newcommand\contentsname{Tabla de contenidos}
\fi
\ifdefined\listfigurename
  \renewcommand*\listfigurename{Lista de Figuras}
\else
  \newcommand\listfigurename{Lista de Figuras}
\fi
\ifdefined\listtablename
  \renewcommand*\listtablename{Lista de Tablas}
\else
  \newcommand\listtablename{Lista de Tablas}
\fi
\ifdefined\figurename
  \renewcommand*\figurename{Figura}
\else
  \newcommand\figurename{Figura}
\fi
\ifdefined\tablename
  \renewcommand*\tablename{Tabla}
\else
  \newcommand\tablename{Tabla}
\fi
}
\@ifpackageloaded{float}{}{\usepackage{float}}
\floatstyle{ruled}
\@ifundefined{c@chapter}{\newfloat{codelisting}{h}{lop}}{\newfloat{codelisting}{h}{lop}[chapter]}
\floatname{codelisting}{Listado}
\newcommand*\listoflistings{\listof{codelisting}{Lista de Listados}}
\makeatother
\makeatletter
\makeatother
\makeatletter
\@ifpackageloaded{caption}{}{\usepackage{caption}}
\@ifpackageloaded{subcaption}{}{\usepackage{subcaption}}
\makeatother
\makeatletter
\@ifpackageloaded{fontawesome5}{}{\usepackage{fontawesome5}}
\makeatother

% From https://tex.stackexchange.com/a/645996/211326
%%% apa7 doesn't want to add appendix section titles in the toc
%%% let's make it do it
\makeatletter
\xpatchcmd{\appendix}
  {\par}
  {\addcontentsline{toc}{section}{\@currentlabelname}\par}
  {}{}
\makeatother

%% Disable longtable counter
%% https://tex.stackexchange.com/a/248395/211326

\usepackage{etoolbox}

\makeatletter
\patchcmd{\LT@caption}
  {\bgroup}
  {\bgroup\global\LTpatch@captiontrue}
  {}{}
\patchcmd{\longtable}
  {\par}
  {\par\global\LTpatch@captionfalse}
  {}{}
\apptocmd{\endlongtable}
  {\ifLTpatch@caption\else\addtocounter{table}{-1}\fi}
  {}{}
\newif\ifLTpatch@caption
\makeatother

\begin{document}

\maketitle


\hypertarget{toc}{}
\tableofcontents
\newpage
\section[Introduction]{Modificación de cadenas}

\setcounter{secnumdepth}{3}

\setlength\LTleft{0pt}




\subsection{🎯 Motivación}\label{motivaciuxf3n}

Los strings son uno de los tipos de datos más manipulados en
programación. A menudo necesitas limpiar texto (eliminar espacios),
cambiar formato (mayúsculas/minúsculas), reemplazar contenido o dividir
texto en partes. Python ofrece métodos integrados que hacen estas
operaciones simples y eficientes sin necesidad de escribir código
complejo.

\subsection{💡 Intuición}\label{intuiciuxf3n}

Imagina que tienes un texto en un procesador de palabras y quieres
aplicar diferentes transformaciones: poner todo en mayúsculas, eliminar
espacios extra, o buscar y reemplazar palabras. Los métodos de string en
Python son como los botones de formato en tu editor de texto: hacen
exactamente eso, pero desde tu código.

\subsection{🧠 Modelo mental}\label{modelo-mental}

\begin{itemize}
\tightlist
\item
  \textbf{Strings son inmutables}: estos métodos \textbf{no modifican}
  el string original, devuelven uno \textbf{nuevo}.
\item
  \textbf{Métodos encadenables}: puedes aplicar múltiples métodos en
  secuencia.
\item
  \textbf{La mayoría devuelven un nuevo string} o una lista (en el caso
  de \texttt{split()}).
\item
  \textbf{Sintaxis}: \texttt{string.metodo()} donde \texttt{string} es
  tu variable o literal.
\end{itemize}

\subsection{📐 Desarrollo}\label{desarrollo}

\subsubsection{1. Cambiar
mayúsculas/minúsculas}\label{cambiar-mayuxfasculasminuxfasculas}

\begin{Shaded}
\begin{Highlighting}[]
\NormalTok{texto }\OperatorTok{=} \StringTok{"Hola Mundo"}

\CommentTok{\# A mayúsculas}
\BuiltInTok{print}\NormalTok{(texto.upper())   }\CommentTok{\# "HOLA MUNDO"}

\CommentTok{\# A minúsculas}
\BuiltInTok{print}\NormalTok{(texto.lower())   }\CommentTok{\# "hola mundo"}

\CommentTok{\# Capitalizar (primera letra mayúscula)}
\BuiltInTok{print}\NormalTok{(texto.capitalize())  }\CommentTok{\# "Hola mundo"}

\CommentTok{\# Título (cada palabra con mayúscula inicial)}
\BuiltInTok{print}\NormalTok{(texto.title())   }\CommentTok{\# "Hola Mundo"}

\CommentTok{\# Swap case (invierte mayúsculas/minúsculas)}
\BuiltInTok{print}\NormalTok{(}\StringTok{"Hola Mundo"}\NormalTok{.swapcase())  }\CommentTok{\# "hOLA mUNDO"}
\end{Highlighting}
\end{Shaded}

\subsubsection{2. Eliminar espacios
(whitespace)}\label{eliminar-espacios-whitespace}

\begin{Shaded}
\begin{Highlighting}[]
\NormalTok{texto }\OperatorTok{=} \StringTok{"  Hola Mundo  "}

\CommentTok{\# Eliminar espacios al inicio y final}
\BuiltInTok{print}\NormalTok{(texto.strip())      }\CommentTok{\# "Hola Mundo"}

\CommentTok{\# Solo al inicio (left)}
\BuiltInTok{print}\NormalTok{(texto.lstrip())     }\CommentTok{\# "Hola Mundo  "}

\CommentTok{\# Solo al final (right)}
\BuiltInTok{print}\NormalTok{(texto.rstrip())     }\CommentTok{\# "  Hola Mundo"}

\CommentTok{\# También funciona con caracteres específicos}
\NormalTok{texto2 }\OperatorTok{=} \StringTok{"***Hola***"}
\BuiltInTok{print}\NormalTok{(texto2.strip(}\StringTok{"*"}\NormalTok{))  }\CommentTok{\# "Hola"}
\end{Highlighting}
\end{Shaded}

\subsubsection{3. Reemplazar texto}\label{reemplazar-texto}

\begin{Shaded}
\begin{Highlighting}[]
\NormalTok{texto }\OperatorTok{=} \StringTok{"Hola Mundo"}

\CommentTok{\# Reemplazar una subcadena}
\BuiltInTok{print}\NormalTok{(texto.replace(}\StringTok{"Mundo"}\NormalTok{, }\StringTok{"Python"}\NormalTok{))  }\CommentTok{\# "Hola Python"}

\CommentTok{\# Reemplazar un carácter}
\BuiltInTok{print}\NormalTok{(texto.replace(}\StringTok{"o"}\NormalTok{, }\StringTok{"0"}\NormalTok{))  }\CommentTok{\# "H0la Mund0"}

\CommentTok{\# Reemplazar múltiples veces (con límite)}
\NormalTok{texto2 }\OperatorTok{=} \StringTok{"uno dos tres cuatro"}
\BuiltInTok{print}\NormalTok{(texto2.replace(}\StringTok{" "}\NormalTok{, }\StringTok{"{-}"}\NormalTok{, }\DecValTok{2}\NormalTok{))  }\CommentTok{\# "uno{-}dos{-}tres cuatro"}
\end{Highlighting}
\end{Shaded}

\subsubsection{4. Dividir texto (split)}\label{dividir-texto-split}

\begin{Shaded}
\begin{Highlighting}[]
\NormalTok{texto }\OperatorTok{=} \StringTok{"Hola Mundo Python"}

\CommentTok{\# Dividir por espacios (por defecto)}
\BuiltInTok{print}\NormalTok{(texto.split())       }\CommentTok{\# [\textquotesingle{}Hola\textquotesingle{}, \textquotesingle{}Mundo\textquotesingle{}, \textquotesingle{}Python\textquotesingle{}]}

\CommentTok{\# Dividir por un separador específico}
\NormalTok{csv }\OperatorTok{=} \StringTok{"manzana,pera,banana"}
\BuiltInTok{print}\NormalTok{(csv.split(}\StringTok{","}\NormalTok{))      }\CommentTok{\# [\textquotesingle{}manzana\textquotesingle{}, \textquotesingle{}pera\textquotesingle{}, \textquotesingle{}banana\textquotesingle{}]}

\CommentTok{\# Dividir con límite}
\BuiltInTok{print}\NormalTok{(csv.split(}\StringTok{","}\NormalTok{, }\DecValTok{1}\NormalTok{))   }\CommentTok{\# [\textquotesingle{}manzana\textquotesingle{}, \textquotesingle{}pera,banana\textquotesingle{}]}

\CommentTok{\# Dividir por saltos de línea}
\NormalTok{multilinea }\OperatorTok{=} \StringTok{"Línea 1}\CharTok{\textbackslash{}n}\StringTok{Línea 2}\CharTok{\textbackslash{}n}\StringTok{Línea 3"}
\BuiltInTok{print}\NormalTok{(multilinea.split(}\StringTok{"}\CharTok{\textbackslash{}n}\StringTok{"}\NormalTok{))  }\CommentTok{\# [\textquotesingle{}Línea 1\textquotesingle{}, \textquotesingle{}Línea 2\textquotesingle{}, \textquotesingle{}Línea 3\textquotesingle{}]}
\end{Highlighting}
\end{Shaded}

\subsubsection{5. Unir texto (join) - el inverso de
split}\label{unir-texto-join---el-inverso-de-split}

\begin{Shaded}
\begin{Highlighting}[]
\CommentTok{\# Unir una lista en un string}
\NormalTok{palabras }\OperatorTok{=}\NormalTok{ [}\StringTok{"Hola"}\NormalTok{, }\StringTok{"Mundo"}\NormalTok{, }\StringTok{"Python"}\NormalTok{]}
\BuiltInTok{print}\NormalTok{(}\StringTok{" "}\NormalTok{.join(palabras))     }\CommentTok{\# "Hola Mundo Python"}
\BuiltInTok{print}\NormalTok{(}\StringTok{"{-}"}\NormalTok{.join(palabras))     }\CommentTok{\# "Hola{-}Mundo{-}Python"}
\BuiltInTok{print}\NormalTok{(}\StringTok{" | "}\NormalTok{.join(palabras))   }\CommentTok{\# "Hola | Mundo | Python"}
\end{Highlighting}
\end{Shaded}

\subsection{🔬 Ejemplos}\label{ejemplos}

\textbf{Mínimo ---} Métodos básicos:

\begin{Shaded}
\begin{Highlighting}[]
\NormalTok{texto }\OperatorTok{=} \StringTok{"  python es genial  "}
\BuiltInTok{print}\NormalTok{(texto.strip().upper())   }\CommentTok{\# "PYTHON ES GENIAL"}
\BuiltInTok{print}\NormalTok{(texto.title())           }\CommentTok{\# "  Python Es Genial  "}
\end{Highlighting}
\end{Shaded}

\textbf{Progresivo ---} Limpieza de datos:

\begin{Shaded}
\begin{Highlighting}[]
\CommentTok{\# Datos sucios de entrada}
\NormalTok{nombre }\OperatorTok{=} \StringTok{"  ana  "}
\NormalTok{email }\OperatorTok{=} \StringTok{"ANA@GMAIL.COM"}
\NormalTok{archivo }\OperatorTok{=} \StringTok{"imagen.JPG"}
\NormalTok{comentario }\OperatorTok{=} \StringTok{"  ¡Excelente!   "}

\CommentTok{\# Limpieza}
\NormalTok{nombre\_limpio }\OperatorTok{=}\NormalTok{ nombre.strip().title()  }\CommentTok{\# "Ana"}
\NormalTok{email\_limpio }\OperatorTok{=}\NormalTok{ email.lower()            }\CommentTok{\# "ana@gmail.com"}
\NormalTok{archivo\_limpio }\OperatorTok{=}\NormalTok{ archivo.lower()        }\CommentTok{\# "imagen.jpg"}
\NormalTok{comentario\_limpio }\OperatorTok{=}\NormalTok{ comentario.strip()  }\CommentTok{\# "¡Excelente!"}

\BuiltInTok{print}\NormalTok{(}\SpecialStringTok{f"Nombre: }\SpecialCharTok{\{}\NormalTok{nombre\_limpio}\SpecialCharTok{\}}\SpecialStringTok{"}\NormalTok{)}
\BuiltInTok{print}\NormalTok{(}\SpecialStringTok{f"Email: }\SpecialCharTok{\{}\NormalTok{email\_limpio}\SpecialCharTok{\}}\SpecialStringTok{"}\NormalTok{)}
\BuiltInTok{print}\NormalTok{(}\SpecialStringTok{f"Archivo: }\SpecialCharTok{\{}\NormalTok{archivo\_limpio}\SpecialCharTok{\}}\SpecialStringTok{"}\NormalTok{)}
\BuiltInTok{print}\NormalTok{(}\SpecialStringTok{f"Comentario: }\SpecialCharTok{\{}\NormalTok{comentario\_limpio}\SpecialCharTok{\}}\SpecialStringTok{"}\NormalTok{)}
\end{Highlighting}
\end{Shaded}

\textbf{Real ---} Procesamiento de CSV y validación:

\begin{Shaded}
\begin{Highlighting}[]
\CommentTok{\# Datos de una línea CSV}
\NormalTok{linea }\OperatorTok{=} \StringTok{"  Ana García, 28, ana@email.com, activo  "}

\CommentTok{\# Limpiar y dividir}
\NormalTok{datos }\OperatorTok{=}\NormalTok{ linea.strip().split(}\StringTok{","}\NormalTok{)}
\NormalTok{nombre }\OperatorTok{=}\NormalTok{ datos[}\DecValTok{0}\NormalTok{].strip().title()}
\NormalTok{edad }\OperatorTok{=} \BuiltInTok{int}\NormalTok{(datos[}\DecValTok{1}\NormalTok{].strip())}
\NormalTok{email }\OperatorTok{=}\NormalTok{ datos[}\DecValTok{2}\NormalTok{].strip().lower()}
\NormalTok{estado }\OperatorTok{=}\NormalTok{ datos[}\DecValTok{3}\NormalTok{].strip().lower()}

\BuiltInTok{print}\NormalTok{(}\SpecialStringTok{f"Nombre: }\SpecialCharTok{\{}\NormalTok{nombre}\SpecialCharTok{\}}\SpecialStringTok{"}\NormalTok{)}
\BuiltInTok{print}\NormalTok{(}\SpecialStringTok{f"Edad: }\SpecialCharTok{\{}\NormalTok{edad}\SpecialCharTok{\}}\SpecialStringTok{"}\NormalTok{)}
\BuiltInTok{print}\NormalTok{(}\SpecialStringTok{f"Email: }\SpecialCharTok{\{}\NormalTok{email}\SpecialCharTok{\}}\SpecialStringTok{"}\NormalTok{)}
\BuiltInTok{print}\NormalTok{(}\SpecialStringTok{f"Estado: }\SpecialCharTok{\{}\NormalTok{estado}\SpecialCharTok{\}}\SpecialStringTok{"}\NormalTok{)}

\CommentTok{\# Validación}
\ControlFlowTok{if} \KeywordTok{not}\NormalTok{ email.endswith(}\StringTok{"@email.com"}\NormalTok{):}
    \BuiltInTok{print}\NormalTok{(}\StringTok{"Email no válido"}\NormalTok{)}

\ControlFlowTok{if}\NormalTok{ estado }\OperatorTok{==} \StringTok{"activo"}\NormalTok{:}
    \BuiltInTok{print}\NormalTok{(}\StringTok{"Usuario activo"}\NormalTok{)}
\end{Highlighting}
\end{Shaded}

\subsection{⚠️ Errores comunes}\label{errores-comunes}

{\def\LTcaptype{none} % do not increment counter
\begin{longtable}[]{@{}
  >{\raggedright\arraybackslash}p{(\linewidth - 4\tabcolsep) * \real{0.3406}}
  >{\raggedright\arraybackslash}p{(\linewidth - 4\tabcolsep) * \real{0.2754}}
  >{\raggedright\arraybackslash}p{(\linewidth - 4\tabcolsep) * \real{0.3841}}@{}}
\toprule\noalign{}
\begin{minipage}[b]{\linewidth}\raggedright
Error
\end{minipage} & \begin{minipage}[b]{\linewidth}\raggedright
Síntoma
\end{minipage} & \begin{minipage}[b]{\linewidth}\raggedright
Solución
\end{minipage} \\
\midrule\noalign{}
\endhead
\bottomrule\noalign{}
\endlastfoot
\textbf{Esperar que el método modifique el original} & El string
original no cambia & Asigna el resultado:
\texttt{texto\ =\ texto.upper()} \\
\textbf{Encadenar métodos en orden incorrecto} &
\texttt{strip().upper()} vs \texttt{upper().strip()} & El orden importa:
primero limpias, luego transformas \\
\textbf{Confundir \texttt{split()} con \texttt{join()}} & Uso incorrecto
de separadores & \texttt{split()} → string a lista; \texttt{join()} →
lista a string \\
\textbf{Llamar método en literal} & \texttt{"texto".upper()} es válido &
Los métodos funcionan en literales también \\
\end{longtable}
}

\subsection{✅ Buenas prácticas}\label{buenas-pruxe1cticas}

\begin{enumerate}
\def\labelenumi{\arabic{enumi}.}
\item
  \textbf{Encadena métodos para limpieza}:

\begin{Shaded}
\begin{Highlighting}[]
\CommentTok{\# ✅ Limpio y eficiente}
\NormalTok{texto\_limpio }\OperatorTok{=}\NormalTok{ texto.strip().lower().replace(}\StringTok{" "}\NormalTok{, }\StringTok{"\_"}\NormalTok{)}

\CommentTok{\# ❌ Más verbose}
\NormalTok{texto }\OperatorTok{=}\NormalTok{ texto.strip()}
\NormalTok{texto }\OperatorTok{=}\NormalTok{ texto.lower()}
\NormalTok{texto }\OperatorTok{=}\NormalTok{ texto.replace(}\StringTok{" "}\NormalTok{, }\StringTok{"\_"}\NormalTok{)}
\end{Highlighting}
\end{Shaded}
\item
  \textbf{Usa \texttt{.strip()} para limpiar entrada de usuario}:

\begin{Shaded}
\begin{Highlighting}[]
\NormalTok{nombre }\OperatorTok{=} \BuiltInTok{input}\NormalTok{(}\StringTok{"Nombre: "}\NormalTok{).strip().title()}
\NormalTok{email }\OperatorTok{=} \BuiltInTok{input}\NormalTok{(}\StringTok{"Email: "}\NormalTok{).strip().lower()}
\end{Highlighting}
\end{Shaded}
\item
  \textbf{Verifica contenido antes de operar}:

\begin{Shaded}
\begin{Highlighting}[]
\ControlFlowTok{if}\NormalTok{ nombre:  }\CommentTok{\# No está vacío}
\NormalTok{    nombre }\OperatorTok{=}\NormalTok{ nombre.strip()}
\end{Highlighting}
\end{Shaded}
\item
  \textbf{Usa \texttt{split()} con límite para casos específicos}:

\begin{Shaded}
\begin{Highlighting}[]
\CommentTok{\# Primera parte antes del primer espacio}
\NormalTok{primero, resto }\OperatorTok{=} \StringTok{"Hola mundo Python"}\NormalTok{.split(}\StringTok{" "}\NormalTok{, }\DecValTok{1}\NormalTok{)}
\BuiltInTok{print}\NormalTok{(primero)  }\CommentTok{\# "Hola"}
\BuiltInTok{print}\NormalTok{(resto)    }\CommentTok{\# "mundo Python"}
\end{Highlighting}
\end{Shaded}
\item
  \textbf{Prefiere \texttt{join()} en lugar de \texttt{+} para unir
  muchos strings}:

\begin{Shaded}
\begin{Highlighting}[]
\CommentTok{\# ✅ Eficiente}
\NormalTok{resultado }\OperatorTok{=} \StringTok{""}\NormalTok{.join(lista\_de\_strings)}

\CommentTok{\# ❌ Ineficiente (crea muchos objetos)}
\NormalTok{resultado }\OperatorTok{=} \StringTok{""}
\ControlFlowTok{for}\NormalTok{ s }\KeywordTok{in}\NormalTok{ lista\_de\_strings:}
\NormalTok{    resultado }\OperatorTok{+=}\NormalTok{ s}
\end{Highlighting}
\end{Shaded}
\end{enumerate}

\subsection{📝 Ejercicio}\label{ejercicio}

\textbf{Intermedio} --- Procesamiento de texto de archivo

Escribe un programa que procese el siguiente texto:

\begin{Shaded}
\begin{Highlighting}[]
\NormalTok{texto }\OperatorTok{=} \StringTok{"""}
\StringTok{  nombre: Ana García  }
\StringTok{  email: ANA@GMAIL.COM  }
\StringTok{  ciudad:  Madrid   }
\StringTok{  profesión:  Ingeniera de Software  }
\StringTok{  años\_exp:  8  }
\StringTok{"""}
\end{Highlighting}
\end{Shaded}

Requisitos:

\begin{enumerate}
\def\labelenumi{\arabic{enumi}.}
\tightlist
\item
  Limpiar cada línea (eliminar espacios)
\item
  Separar cada línea en clave: valor usando \texttt{:}
\item
  Limpiar las claves (strip, lower)
\item
  Limpiar los valores (strip, formatear según el caso)
\item
  Mostrar los datos formateados
\end{enumerate}

\textbf{Pista:} usa \texttt{split("\textbackslash{}n")} para dividir el
texto en líneas.

\begin{center}\rule{0.5\linewidth}{0.5pt}\end{center}

\textbf{Solución (no mires hasta haberlo intentado):}

\begin{Shaded}
\begin{Highlighting}[]
\NormalTok{texto }\OperatorTok{=} \StringTok{"""}
\StringTok{  nombre: Ana García  }
\StringTok{  email: ANA@GMAIL.COM  }
\StringTok{  ciudad:  Madrid   }
\StringTok{  profesión:  Ingeniera de Software  }
\StringTok{  años\_exp:  8  }
\StringTok{"""}

\CommentTok{\# Dividir en líneas y limpiar}
\NormalTok{lineas }\OperatorTok{=}\NormalTok{ texto.strip().split(}\StringTok{"}\CharTok{\textbackslash{}n}\StringTok{"}\NormalTok{)}
\NormalTok{datos }\OperatorTok{=}\NormalTok{ \{\}}

\ControlFlowTok{for}\NormalTok{ linea }\KeywordTok{in}\NormalTok{ lineas:}
    \ControlFlowTok{if} \StringTok{":"} \KeywordTok{in}\NormalTok{ linea:}
        \CommentTok{\# Dividir por el primer ":"}
\NormalTok{        clave, valor }\OperatorTok{=}\NormalTok{ linea.split(}\StringTok{":"}\NormalTok{, }\DecValTok{1}\NormalTok{)}

        \CommentTok{\# Limpiar}
\NormalTok{        clave }\OperatorTok{=}\NormalTok{ clave.strip().lower()}
\NormalTok{        valor }\OperatorTok{=}\NormalTok{ valor.strip()}

        \CommentTok{\# Formatear según el tipo de dato}
        \ControlFlowTok{if}\NormalTok{ clave }\OperatorTok{==} \StringTok{"email"}\NormalTok{:}
\NormalTok{            valor }\OperatorTok{=}\NormalTok{ valor.lower()}
        \ControlFlowTok{elif}\NormalTok{ clave }\OperatorTok{==} \StringTok{"nombre"}\NormalTok{:}
\NormalTok{            valor }\OperatorTok{=}\NormalTok{ valor.title()}
        \ControlFlowTok{elif}\NormalTok{ clave }\OperatorTok{==} \StringTok{"ciudad"}\NormalTok{:}
\NormalTok{            valor }\OperatorTok{=}\NormalTok{ valor.strip().title()}
        \ControlFlowTok{elif}\NormalTok{ clave }\OperatorTok{==} \StringTok{"años\_exp"}\NormalTok{:}
\NormalTok{            valor }\OperatorTok{=} \BuiltInTok{int}\NormalTok{(valor)  }\CommentTok{\# Convertir a entero}

\NormalTok{        datos[clave] }\OperatorTok{=}\NormalTok{ valor}

\CommentTok{\# Mostrar resultados}
\BuiltInTok{print}\NormalTok{(}\StringTok{"=== Datos Procesados ==="}\NormalTok{)}
\ControlFlowTok{for}\NormalTok{ clave, valor }\KeywordTok{in}\NormalTok{ datos.items():}
    \BuiltInTok{print}\NormalTok{(}\SpecialStringTok{f"}\SpecialCharTok{\{}\NormalTok{clave}\SpecialCharTok{\}}\SpecialStringTok{: }\SpecialCharTok{\{}\NormalTok{valor}\SpecialCharTok{\}}\SpecialStringTok{"}\NormalTok{)}

\BuiltInTok{print}\NormalTok{(}\SpecialStringTok{f"}\CharTok{\textbackslash{}n}\SpecialStringTok{Tipo de años\_exp: }\SpecialCharTok{\{}\BuiltInTok{type}\NormalTok{(datos[}\StringTok{\textquotesingle{}años\_exp\textquotesingle{}}\NormalTok{])}\SpecialCharTok{\}}\SpecialStringTok{"}\NormalTok{)}
\end{Highlighting}
\end{Shaded}

\subsection{🗂️ Síntesis}\label{suxedntesis}

\begin{itemize}
\tightlist
\item
  \textbf{\texttt{.upper()} / \texttt{.lower()}}: cambio de
  mayúsculas/minúsculas.
\item
  \textbf{\texttt{.strip()} / \texttt{.lstrip()} / \texttt{.rstrip()}}:
  eliminar espacios al inicio/final.
\item
  \textbf{\texttt{.replace()}}: reemplaza texto por otro.
\item
  \textbf{\texttt{.split()}}: divide un string en lista según separador.
\item
  \textbf{\texttt{.join()}}: une una lista en string (inverso de split).
\item
  \textbf{Inmutabilidad}: todos los métodos devuelven strings nuevos.
\item
  \textbf{Encadenamiento}: puedes aplicar múltiples métodos
  secuencialmente.
\item
  \textbf{\texttt{title()} / \texttt{capitalize()}}: formatos especiales
  de capitalización.
\end{itemize}

\subsection{🔗 Qué sigue}\label{quuxe9-sigue}

\begin{itemize}
\tightlist
\item
  \textbf{Relacionado:} métodos de string avanzados, f-strings,
  expresiones regulares
\item
  \textbf{Siguiente tema:} Concatenación y formateo de strings ---
  aprende las mejores formas de combinar texto y variables.
\end{itemize}

\section{Publicaciones Similares}\label{publicaciones-similares}

Si te interesó este artículo, te recomendamos que explores otros blogs y
recursos relacionados que pueden ampliar tus conocimientos. Aquí te dejo
algunas sugerencias:

\begin{enumerate}
\def\labelenumi{\arabic{enumi}.}
\tightlist
\item
  \href{https://numerus-scriptum.netlify.app/python/2020-06-19-instalacion-de-anaconda/index.pdf}{\faIcon{file-pdf}}
  \href{https://numerus-scriptum.netlify.app/python/2020-06-19-instalacion-de-anaconda}{Instalacion
  De Anaconda}
\item
  \href{https://numerus-scriptum.netlify.app/python/2020-06-20-configurar-entorno-virtual-python-anaconda/index.pdf}{\faIcon{file-pdf}}
  \href{https://numerus-scriptum.netlify.app/python/2020-06-20-configurar-entorno-virtual-python-anaconda}{Configurar
  Entorno Virtual Python Anaconda}
\item
  \href{https://numerus-scriptum.netlify.app/python/2021-04-17-01-introducion-a-la-programacion-con-python/index.pdf}{\faIcon{file-pdf}}
  \href{https://numerus-scriptum.netlify.app/python/2021-04-17-01-introducion-a-la-programacion-con-python}{01
  Introducion A La Programacion Con Python}
\item
  \href{https://numerus-scriptum.netlify.app/python/2021-05-31-02-variables-expresiones-y-statements-con-python/index.pdf}{\faIcon{file-pdf}}
  \href{https://numerus-scriptum.netlify.app/python/2021-05-31-02-variables-expresiones-y-statements-con-python}{02
  Variables Expresiones Y Statements Con Python}
\item
  \href{https://numerus-scriptum.netlify.app/python/2021-06-07-03-objetos-de-python/index.pdf}{\faIcon{file-pdf}}
  \href{https://numerus-scriptum.netlify.app/python/2021-06-07-03-objetos-de-python}{03
  Objetos De Python}
\item
  \href{https://numerus-scriptum.netlify.app/python/2021-06-14-04-ejecucion-condicional-con-python/index.pdf}{\faIcon{file-pdf}}
  \href{https://numerus-scriptum.netlify.app/python/2021-06-14-04-ejecucion-condicional-con-python}{04
  Ejecucion Condicional Con Python}
\item
  \href{https://numerus-scriptum.netlify.app/python/2021-06-21-05-iteraciones-con-python/index.pdf}{\faIcon{file-pdf}}
  \href{https://numerus-scriptum.netlify.app/python/2021-06-21-05-iteraciones-con-python}{05
  Iteraciones Con Python}
\item
  \href{https://numerus-scriptum.netlify.app/python/2021-08-16-06-funciones-con-python/index.pdf}{\faIcon{file-pdf}}
  \href{https://numerus-scriptum.netlify.app/python/2021-08-16-06-funciones-con-python}{06
  Funciones Con Python}
\item
  \href{https://numerus-scriptum.netlify.app/python/2021-08-23-07-dataframes-con-python/index.pdf}{\faIcon{file-pdf}}
  \href{https://numerus-scriptum.netlify.app/python/2021-08-23-07-dataframes-con-python}{07
  Dataframes Con Python}
\item
  \href{https://numerus-scriptum.netlify.app/python/2021-11-29-08-prediccion-y-metrica-de-performance-con-python/index.pdf}{\faIcon{file-pdf}}
  \href{https://numerus-scriptum.netlify.app/python/2021-11-29-08-prediccion-y-metrica-de-performance-con-python}{08
  Prediccion Y Metrica De Performance Con Python}
\item
  \href{https://numerus-scriptum.netlify.app/python/2021-12-06-09-metodos-de-machine-learning-para-clasificacion-con-python/index.pdf}{\faIcon{file-pdf}}
  \href{https://numerus-scriptum.netlify.app/python/2021-12-06-09-metodos-de-machine-learning-para-clasificacion-con-python}{09
  Metodos De Machine Learning Para Clasificacion Con Python}
\item
  \href{https://numerus-scriptum.netlify.app/python/2021-12-13-10-metodos-de-machine-learning-para-regresion-con-python/index.pdf}{\faIcon{file-pdf}}
  \href{https://numerus-scriptum.netlify.app/python/2021-12-13-10-metodos-de-machine-learning-para-regresion-con-python}{10
  Metodos De Machine Learning Para Regresion Con Python}
\item
  \href{https://numerus-scriptum.netlify.app/python/2025-05-10-visualizacion-de-datos-con-python/index.pdf}{\faIcon{file-pdf}}
  \href{https://numerus-scriptum.netlify.app/python/2025-05-10-visualizacion-de-datos-con-python}{Visualizacion
  De Datos Con Python}
\end{enumerate}

Esperamos que encuentres estas publicaciones igualmente interesantes y
útiles. ¡Disfruta de la lectura!






\end{document}
